% 02_modelo.tex — Modelo matemático (ILP)
\section{Modelo matemático}
\label{sec:modelo}

\subsection{Conjuntos y parámetros}
\label{sec:modelo:conjuntos}

Los conjuntos $X$, $Y$, $Z$ y la colección $T$ se definen como en la
Sección~\ref{sec:problema}.  Para cada tripleta $t\in T$ y cada elemento
$i\in X$, $j\in Y$, $k\in Z$ definimos los \emph{indicadores de incidencia}:
\[
a^{X}_{t,i} = \begin{cases} 1 & \text{si } i = \pi_X(t) \\ 0 & \text{en otro caso} \end{cases},
\qquad
a^{Y}_{t,j} = \begin{cases} 1 & \text{si } j = \pi_Y(t) \\ 0 & \text{en otro caso} \end{cases},
\qquad
a^{Z}_{t,k} = \begin{cases} 1 & \text{si } k = \pi_Z(t) \\ 0 & \text{en otro caso} \end{cases}.
\]
Como cada tripleta tiene exactamente un elemento por dimensión,
$\sum_{i\in X} a^{X}_{t,i}=1$ para todo $t\in T$ (análogo para $Y$, $Z$).

\subsection{Variables de decisión}
\label{sec:modelo:variables}

Se define una variable binaria por tripleta:
\[
s_t \in \{0,1\},\quad t\in T,
\qquad
s_t = 1 \;\Longleftrightarrow\; t\in M.
\]
La notación $s_t$ (\emph{selector}) evita colisión con los elementos
$x\in X$, que ya usan la letra~$x$.  Total: $m$ variables binarias.

\subsection{Función objetivo}
\label{sec:modelo:objetivo}

Maximizar la cardinalidad del matching:
\[
\max \;\sum_{t\in T} s_t.
\]

\subsection{Restricciones}
\label{sec:modelo:restricciones}

Cada elemento puede ser cubierto por a lo más una tripleta.  Esto da lugar a
tres familias de $n$ restricciones cada una:

\begin{align}
\text{(C-X)} &\quad \sum_{t\in T} a^{X}_{t,i}\,s_t \;\le\; 1, \qquad \forall\,i\in X, \label{eq:cx}\\
\text{(C-Y)} &\quad \sum_{t\in T} a^{Y}_{t,j}\,s_t \;\le\; 1, \qquad \forall\,j\in Y, \label{eq:cy}\\
\text{(C-Z)} &\quad \sum_{t\in T} a^{Z}_{t,k}\,s_t \;\le\; 1, \qquad \forall\,k\in Z. \label{eq:cz}
\end{align}

Total: $3n$ restricciones lineales más $m$ restricciones de integralidad.

\subsection{Modelo completo}
\label{sec:modelo:completo}

\begin{equation}
\label{eq:ilp}
\begin{array}{rl}
\displaystyle\max & \displaystyle\sum_{t\in T} s_t \\[1.2ex]
\text{s.a.} & \displaystyle\sum_{t\ni i} s_t \;\le\; 1, \quad \forall\,i\in X, \\[1.0ex]
            & \displaystyle\sum_{t\ni j} s_t \;\le\; 1, \quad \forall\,j\in Y, \\[1.0ex]
            & \displaystyle\sum_{t\ni k} s_t \;\le\; 1, \quad \forall\,k\in Z, \\[1.0ex]
            & s_t \in \{0,1\}, \quad \forall\,t\in T,
\end{array}
\tag{P$_\text{ILP}$}
\end{equation}
donde la notación $t\ni i$ abrevia $a^{X}_{t,i}=1$.

\subsection{Equivalencia decisión \texorpdfstring{$\Leftrightarrow$}{<->} optimización}
\label{sec:modelo:equivalencia}

Sea $\mathcal{D}(I)\in\{\text{Sí},\text{No}\}$ el oráculo de decisión.
La equivalencia es polinomial en ambos sentidos:
\[
\mathcal{D}(I) = \text{Sí}
\;\Longleftrightarrow\;
\OPT(I) = n.
\]
\noindent
\textbf{Decisión $\Rightarrow$ Optimización (padding).}
Para responder ``$\OPT(I)\ge k$?'' construimos $I_k$ aumentando $I$ con
$n-k$ tripletas \emph{dummy} completamente disjuntas.
Un matching perfecto en $I_k$ existe si y sólo si $\OPT(I)\ge k$.
Con búsqueda binaria, $O(\log n)$ llamadas al oráculo de decisión determinan
$\OPT(I)$; la construcción de cada $I_k$ es lineal.

\subsection{Formulación como Exact Cover tripartito}
\label{sec:modelo:x3c}

Toda instancia $I=(X,Y,Z,T)$ de 3DM induce una instancia de
\textit{Exact Cover by 3-Sets} (X3C)~\cite{garey1979}:
\[
U = X\cup Y\cup Z,\qquad |U|=3n,\qquad
\mathcal{C} = \bigl\{\{x,y,z\} : (x,y,z)\in T\bigr\}.
\]
Un \emph{exact cover} $\mathcal{C}^*\subseteq\mathcal{C}$ con $|\mathcal{C}^*|=n$
corresponde biyectivamente a un matching perfecto en $I$.
Esta correspondencia justifica los algoritmos de
\emph{exact cover} (tipo Dancing Links) como referencia teórica; en este
trabajo, la lógica equivalente se implementa a mano mediante backtracking
(véase Sección~\ref{sec:implementacion}).

\subsection{Relajación LP y brecha de integralidad}
\label{sec:modelo:lp}

A diferencia del matching bipartito clásico, cuya matriz de restricciones es
\emph{totalmente unimodular} y garantiza soluciones enteras para la relajación
LP~\cite{garey1979}, la matriz de incidencia tripleta--elemento de 3DM
\textbf{no} es totalmente unimodular.

\begin{lemma}[Gap LP--ILP]
\label{lem:gap}
Existen instancias de 3DM en que la solución óptima de la relajación LP es
estrictamente mayor que $\OPT(I)$.
\end{lemma}
\begin{proof}
Sea $n=2$ y $T=\{t_1{=}(x_1,y_1,z_1),\;t_2{=}(x_1,y_2,z_2),\;t_3{=}(x_2,y_1,z_2)\}$.
Cada par de tripletas comparte algún elemento, por lo que el óptimo entero es
$\OPT_{\text{ILP}}=1$.
La asignación $s_1=s_2=s_3=\tfrac{1}{2}$ es factible en la relajación LP
(cada restricción suma a lo más 1) y da valor $\OPT_{\text{LP}}=\tfrac{3}{2}$.
La brecha $\OPT_{\text{LP}}-\OPT_{\text{ILP}}=\tfrac{1}{2}$ descarta la
unimodularidad total.\qed
\end{proof}

Este resultado justifica recurrir a \emph{backtracking} con poda en lugar de
un simple redondeo de la relajación LP.

\subsection{Resumen del modelo}
\label{sec:modelo:resumen}

\begin{table}[htbp]
\centering
\caption{Características del modelo ILP para max-3DM.}
\label{tab:modelo_resumen}
\begin{tabular}{ll}
\toprule
Característica & Valor \\
\midrule
Variables & $m$ binarias ($s_t\in\{0,1\}$) \\
Restricciones lineales & $3n$ \\
Sentido objetivo & maximización \\
Tipo & ILP (programación entera 0--1) \\
Relajación LP entera & no (en general) \\
Cota dual trivial & $\OPT(I)\le n$ \\
Tamaño peor caso & $m=O(n^3)$ \\
\bottomrule
\end{tabular}
\end{table}
